%!TEX root = ../main.tex
% GLOSARIO: términos que necesitan definición dentro del documento.
%
% Definir:   \newglossaryentry{clave}{name={Término}, description={Definición}}
% Usar:      \gls{clave}      -> Término (con enlace al glosario)
%            \Gls{clave}      -> Término (con mayúscula inicial)
%            \glspl{clave}    -> plural
% Solo aparecen en el glosario los términos que se usan en el texto.
% Las siglas van en siglas.tex y los símbolos en simbolos.tex.

\newglossaryentry{g:mvp}{
    name={producto mínimo viable},
    description={Versión funcional de un producto con el conjunto mínimo
    de características que entrega valor real a los usuarios y permite
    aprender de su uso. A diferencia de un prototipo, no simula: funciona}
}

\newglossaryentry{prototipo}{
    name={prototipo},
    description={Versión preliminar de un sistema construida para validar
    una idea, un flujo o una interfaz. Puede simular partes del
    comportamiento y no está pensada para operar en producción}
}

\newglossaryentry{stakeholder}{
    name={parte interesada},
    description={Persona, grupo u organización que puede afectar o verse
    afectada por el proyecto: usuarios, clientes, patrocinadores,
    reguladores y el propio equipo}
}

\newglossaryentry{exposicion}{
    name={exposición al riesgo},
    description={Producto de la probabilidad de ocurrencia de un riesgo por
    su impacto. Se usa para ordenar la matriz de riesgos y decidir cuáles
    se tratan primero}
}

\newglossaryentry{g:van}{
    name={valor actual neto},
    description={Suma de los flujos de caja futuros de un proyecto
    descontados a una tasa que refleja su riesgo, menos la inversión
    inicial. Un valor positivo indica que el proyecto crea valor}
}

\newglossaryentry{g:marketplace}{
    name={marketplace},
    description={Plataforma digital que conecta a dos grupos de usuarios
    distintos --- en este proyecto, clientes que buscan contratar
    servicios de eventos y salones/proveedores que ofrecen esos
    servicios--- y les permite comparar, comunicarse y contratar entre sí}
}

\newglossaryentry{g:sistema-recomendacion}{
    name={sistema de recomendación},
    description={Herramienta que procesa información sobre preferencias,
    comportamiento y características de los usuarios para sugerir
    automáticamente las opciones más adecuadas}
}

\newglossaryentry{g:gemelo-digital}{
    name={Gemelo Digital},
    plural={Gemelos Digitales},
    description={Representación virtual de un objeto, proceso o sistema
    real que se actualiza continuamente con nuevos datos, permitiendo
    simular distintos escenarios a medida que cambia la información del
    sistema real. A diferencia de un modelo estático o una simulación
    tradicional, evoluciona junto con el sistema que representa}
}

\newglossaryentry{g:gestion-integral-eventos}{
    name={gestión integral de eventos},
    description={Conjunto de actividades relacionadas con la
    planificación, coordinación y administración de todos los recursos y
    servicios involucrados en la realización de un evento}
}

\newglossaryentry{g:economia-plataformas}{
    name={economía de plataformas},
    description={Modelo de negocio digital en el que una plataforma crea
    valor principalmente al facilitar las interacciones entre dos o más
    grupos de usuarios interdependientes, también llamado mercado de dos
    lados \citep{rochettirole2006}}
}

\newglossaryentry{g:big-data}{
    name={Big Data},
    description={Conjuntos de datos de gran volumen, velocidad y
    variedad, base sobre la que operan los sistemas de recomendación y
    los Gemelos Digitales}
}

\newglossaryentry{g:cloud-computing}{
    name={computación en la nube},
    description={Modelo de provisión de infraestructura, almacenamiento y
    procesamiento a través de Internet, sin administrar servidores
    físicos propios}
}
